% ======================= Section 3: Waveforms beyond OFDM (S12-S20) =======
\section{Waveforms Beyond OFDM}

% --- S12 -----------------------------------------------------------------
\begin{frame}{CP-OFDM recap: why it won 4G and 5G}
  \begin{equation*}
    x[n]=\frac{1}{\sqrt N}\sum_{k=0}^{N-1}X_k\,e^{j2\pi kn/N},
    \qquad\text{CP: linear conv.}\;\to\;\text{circular conv.}\;\to\;
    \text{1-tap FDE}
  \end{equation*}
  \begin{itemize}\itemsep2pt
    \item One FFT turns a frequency-selective channel into $N$ flat channels
    \item MIMO-friendly (per-subcarrier processing), flexible numerology,
          gigantic ecosystem
    \item Hidden assumption: \alert{channel constant over one symbol}
          (quasi-static)
    \item That assumption is precisely what high mobility and high carriers
          break
  \end{itemize}
  \note{[S12, 1 min | clock 0:13] Quick refresher --- the audience knows OFDM,
  so frame it as ``what exactly did we sign up for?'' The CP-makes-it-circular
  trick is the whole magic; note it costs 7--25 percent overhead. Set up the
  fall: the one-tap equalizer only exists if the channel sits still for a
  symbol.}
\end{frame}

% --- S13 -----------------------------------------------------------------
\begin{frame}{Where OFDM breaks: high Doppler, high bands}
  \begin{columns}[T]
    \column{0.5\textwidth}
    \begin{equation*}
      P_{\mathrm{ICI}}\;\lesssim\;\frac{(\pi f_d T_s)^{2}}{3}
    \end{equation*}
    \begin{itemize}\itemsep2pt
      \item Doppler spread destroys orthogonality $\Rightarrow$ ICI
      \item 500\,km/h at 28\,GHz: $f_d\!\approx\!13$\,kHz
            ($\varepsilon\!\approx\!0.11$ even at 120\,kHz spacing)
      \item PAPR $\sim$10--12\,dB $\Rightarrow$ PA back-off --- worst fit for
            THz PAs
      \item Phase noise adds CPE + ICI on top
    \end{itemize}
    \column{0.48\textwidth}
    \centering
    \includegraphics[width=\textwidth]{fig_ici_doppler.png}
  \end{columns}
  \note{[S13, 1 min | clock 0:14] Walk the log-log plot: ICI grows with the
  square of normalized Doppler; the pink marker is our 500 km/h working point
  --- ICI around $-17$\,dB, i.e.\ an SINR ceiling near 17 dB no matter how much
  power you pour in. That ceiling becomes the BER floor in Demo 1.}
\end{frame}

% --- S14 -----------------------------------------------------------------
\begin{frame}{DFT-s-OFDM and the single-carrier comeback}
  \begin{columns}[T]
    \column{0.5\textwidth}
    \begin{equation*}
      \mathbf X = F_M\,\mathbf s \quad\text{(DFT spread)}
    \end{equation*}
    \begin{itemize}\itemsep2pt
      \item Quasi single-carrier envelope: $\approx$3\,dB PAPR advantage
      \item Already the 5G uplink workhorse --- proven at scale
      \item Sub-THz proposals: SC-FDE blocks; equalizer complexity moves to RX
      \item Leading candidate for 6G high bands in early 3GPP discussion
    \end{itemize}
    \column{0.48\textwidth}
    \centering
    \includegraphics[width=\textwidth]{fig_papr_ccdf.png}
  \end{columns}
  \note{[S14, 1 min | clock 0:15] The CCDF plot is simulated, not schematic:
  QPSK, 1200 tones. Read off the 1e-3 point: about 3 dB less back-off, which
  at THz translates almost 1:1 into transmit power and hence link range. The
  old single-carrier vs multicarrier war is back, and this time single-carrier
  has a real case.}
\end{frame}

% --- S15 -----------------------------------------------------------------
\begin{frame}{Filtered multicarrier: lessons of the 5G bake-off}
  \begin{itemize}\itemsep2pt
    \item Candidates circa 2016: FBMC, UFMC, GFDM, f-OFDM --- better spectral
          confinement, relaxed synchronicity
    \item Why they lost: MIMO integration pain (FBMC/OQAM), complexity,
          marginal gains once windowing was allowed
    \item What survived: spectral shaping as a \emph{transmitter
          implementation choice} below the standard
    \item \textbf{Meta-lesson:} ecosystem inertia is a design constraint ---
          a radical waveform needs \alert{radical} gains, not 0.5\,dB
  \end{itemize}
  \note{[S15, 1 min | clock 0:16] Flex slide --- compress to 20 seconds if
  behind (say only the meta-lesson). This history is why OTFS/AFDM must clear
  a high bar: they must enable something OFDM fundamentally cannot do, which
  is exactly the high-Doppler regime, not just look prettier in a PSD plot.}
\end{frame}

% --- S16 -----------------------------------------------------------------
\begin{frame}{OTFS I: thinking in delay--Doppler}
  \begin{equation*}
    r(t)=\iint h(\tau,\nu)\,s(t-\tau)\,e^{j2\pi\nu(t-\tau)}\,
    \mathrm d\tau\,\mathrm d\nu
  \end{equation*}
  \begin{columns}[T]
    \column{0.55\textwidth}
    \begin{itemize}\itemsep1pt
      \item In delay--Doppler (DD) the channel is \alert{sparse and
            quasi-static}: a few $(\tau_i,\nu_i)$ taps $=$ physical reflectors
            with range \& velocity
      \item Each DD symbol rides the \emph{entire} TF channel
            $\Rightarrow$ full-diversity potential
    \end{itemize}
    \column{0.43\textwidth}
    \centering
    \includegraphics[width=\textwidth]{fig_tf_vs_dd.png}
  \end{columns}
  \note{[S16, 1.5 min | clock 0:17] The domain-change story, told with our own
  simulation: left panel, the TF response of the exact channel used in Demo 1
  --- fading everywhere; right panel, the SAME channel in DD --- a handful of
  still taps. Analogy: TF is the weather, DD is the map of what causes the
  weather.}
\end{frame}

% --- S17 -----------------------------------------------------------------
\begin{frame}{OTFS II: transceiver structure}
  \begin{equation*}
    X_{\mathrm{tf}}[n,m]=\tfrac{1}{\sqrt{NM}}
    \sum_{k=0}^{N-1}\sum_{l=0}^{M-1}x[k,l]\,
    e^{\,j2\pi\left(\frac{nk}{N}-\frac{ml}{M}\right)}\qquad\text{(ISFFT)}
  \end{equation*}
  \centering
  \resizebox{0.98\textwidth}{!}{% OTFS transceiver chain (slide 17) -- native TikZ
\begin{tikzpicture}[node distance=4mm and 6mm]
  \node[lab] (x) {$x[k,l]$\\[-2pt]\scriptsize DD data};
  \node[blk, right=of x] (isfft) {ISFFT};
  \node[blk, right=of isfft] (heis) {Heisenberg\\transform\\ \scriptsize(OFDM mod.)};
  \node[blkb, right=of heis, fill=cblue!15] (ch) {$h(\tau,\nu)$\\ \scriptsize doubly\\[-3pt]\scriptsize dispersive};
  \node[blk, right=of ch] (wig) {Wigner\\transform\\ \scriptsize(OFDM demod.)};
  \node[blk, right=of wig] (sfft) {SFFT};
  \node[lab, right=of sfft] (y) {$y[k,l]$\\[-2pt]\scriptsize DD out};

  \draw[arr] (x) -- (isfft);
  \draw[arr] (isfft) -- node[lab, above]{$X_{\mathrm{tf}}[n,m]$} (heis);
  \draw[arr] (heis) -- node[lab, above]{$s(t)$} (ch);
  \draw[arr] (ch) -- node[lab, above]{$r(t)$} (wig);
  \draw[arr] (wig) -- node[lab, above]{$Y_{\mathrm{tf}}[n,m]$} (sfft);
  \draw[arr] (sfft) -- (y);

  \node[draw=gray, dashed, rounded corners, fit=(isfft)(heis),
        inner sep=2.5mm, label={[lab]below:2D precoding across the frame}] {};
  \node[draw=gray, dashed, rounded corners, fit=(wig)(sfft),
        inner sep=2.5mm, label={[lab]below:back to delay--Doppler}] {};
\end{tikzpicture}
}
  \begin{itemize}\itemsep1pt
    \item OTFS $=$ 2D pre-/post-processing wrapped around an OFDM modem
          $\Rightarrow$ attractive migration path
  \end{itemize}
  \note{[S17, 1.5 min | clock 0:18:30] Walk the chain left to right: ISFFT
  spreads every DD symbol over the whole TF frame; Heisenberg transform is
  just the OFDM modulator in disguise; receiver undoes both. Emphasize the
  implementation message: no exotic hardware, an extra 2D FFT pair. Pilots: a
  single DD impulse with a guard box sounds the whole channel.}
\end{frame}

% --- S18 -----------------------------------------------------------------
\begin{frame}{OTFS III: input--output relation and detection}
  \begin{equation*}
    y[k,l]\;\approx\;\sum_{i}h_i\;
    x\!\left[(k-k_i)_N,\,(l-l_i)_M\right]e^{j2\pi\phi_i}\;+\;w[k,l]
  \end{equation*}
  \begin{itemize}\itemsep2pt
    \item A \alert{2D twisted circular convolution}: each tap shifts data in
          delay and Doppler and adds a phase
    \item Sparsity $\Rightarrow$ message-passing / low-complexity MMSE
          detectors
    \item Costs: frame-level latency, detection complexity, fractional-Doppler
          leakage
    \item Status check (honest): rich literature, strong results ---
          \emph{not yet} in 3GPP normative work
  \end{itemize}
  \note{[S18, 1 min | clock 0:20] The relation is OFDM's one-tap story
  upgraded to 2D: instead of a scalar per subcarrier, a small set of shifts of
  the whole grid. Detection is where OTFS pays its bill --- mention MP
  detectors exploit the sparsity. Keep the standards status honest; that
  credibility pays off in Q\&A.}
\end{frame}

% --- S19 -----------------------------------------------------------------
\begin{frame}{AFDM: the chirp-based challenger}
  \begin{equation*}
    \phi_m[n]=\tfrac{1}{\sqrt N}\,
    e^{\,j2\pi\left(c_1 n^{2}+\frac{mn}{N}+c_2 m^{2}\right)}
    \qquad\text{(discrete affine Fourier basis)}
  \end{equation*}
  \begin{itemize}\itemsep2pt
    \item Multiplex on \alert{chirps}; tune $c_1$ to the expected Doppler span
          $\Rightarrow$ all paths separate in the DAFT domain
    \item Full diversity in doubly dispersive channels with \emph{1D,
          single-symbol} processing (vs OTFS frame-level)
    \item FFT-based implementation; rising fast in the literature since 2023
    \item Open: MIMO integration, pilot design, robustness studies at scale
  \end{itemize}
  \note{[S19, 1 min | clock 0:21] One-slide take: AFDM gets OTFS-like
  diversity with lower latency because it works symbol-by-symbol. The chirp
  parameter is matched to max Doppler --- a nice example of channel-aware basis
  design. If pressed for time, the sentence to keep: ``same benefit, smaller
  reorganization of the transceiver.''}
\end{frame}

% --- S20 -----------------------------------------------------------------
\begin{frame}{Waveform scorecard}
  \centering
  \includegraphics[width=0.62\textwidth]{fig_waveform_scorecard.png}
  \demoslot{1 (3.5 min)}{\textbf{OFDM vs OTFS at 500\,km/h} --- watch the BER
    floor appear on screen, live. \hfill\src{matlab/demo1\_ofdm\_vs\_otfs}}
  \note{[S20 + DEMO 1, 1+3.5 min | clock 0:22] Summarize: OFDM keeps the
  low-mobility mainstream; DD-domain waveforms own high mobility and pair
  naturally with sensing. THEN SWITCH TO MATLAB (pre-typed command:
  run\_demo). While it runs: recall the ICI ceiling from S13, predict the
  floor, let the plot confirm it. Expected: OFDM floors at BER a few times
  1e-2, OTFS falls through 1e-4. If MATLAB misbehaves: appendix slide A1 has
  the exact expected output. Clock check: demo done by 0:26:30.}
\end{frame}
